% Generated draft: requires mathematical and source review.
\documentclass[11pt]{article}
\usepackage[margin=25mm]{geometry}
\usepackage{fontspec}
\setmainfont{Latin Modern Roman}
\usepackage{amsmath,amssymb,mathtools,unicode-math}
\setmathfont{Latin Modern Math}
\usepackage{enumitem,needspace,xurl,hyperref}
\setlength{\parindent}{0pt}
\setlength{\parskip}{5pt}
\setlength{\emergencystretch}{3em}
\newcommand{\N}{\mathbb{N}}
\newcommand{\Z}{\mathbb{Z}}
\newcommand{\Q}{\mathbb{Q}}
\newcommand{\R}{\mathbb{R}}
\newcommand{\C}{\mathbb{C}}
\newcommand{\F}{\mathbb{F}}
\newcommand{\A}{\mathbb{A}}
\newcommand{\PP}{\mathbb{P}}
\newcommand{\E}{\mathbb{E}}
\newcommand{\eps}{\varepsilon}
\newcommand{\dd}{\,\mathrm{d}}
\begin{document}
\section*{UnsolvedMath ID: 30004716}
\textbf{Source problem code:} catalogue entry
\subsection{Definitions and mathematical statement}
The supplied record poses a research question rather than a fully quantified conjecture. The following notation describes a conventional framework; the manifold hypotheses and the intended classes remain to be specified.

Let \(M\) be a smooth manifold. An initial data set on \(M\) consists of a Riemannian metric \(g\) and a smooth symmetric covariant \(2\)-tensor \(k\). Under the conventional constraint normalization with zero cosmological constant, define the energy density \(\rho\) and momentum density \(J\) by
\[
2\rho(g,k)=\operatorname{Scal}_g+(\operatorname{tr}_g k)^2-|k|_g^2,
\qquad
J(g,k)=\operatorname{div}_g k-\mathrm d(\operatorname{tr}_g k).
\]
The strict dominant energy condition is the pointwise inequality
\[
\rho(g,k)(x)>|J(g,k)(x)|_g\qquad\text{for every }x\in M.
\]
This normalization is an explicit interpretive convention, not an additional hypothesis stated in the source.

Write
\[
\mathcal R^{>0}(M)=\{g:\operatorname{Scal}_g(x)>0\text{ for every }x\in M\},
\]
\[
\mathcal I^{\mathrm{strict}}(M)
=\{(g,k):\rho(g,k)(x)>|J(g,k)(x)|_g\text{ for every }x\in M\}.
\]
Any restrictions on the metrics and tensors, including boundary or asymptotic conditions, must be specified before these sets define the intended spaces. With compatible restrictions and topologies, the time-symmetric assignment gives a map
\[
\iota:\mathcal R^{>0}(M)\longrightarrow\mathcal I^{\mathrm{strict}}(M),
\qquad \iota(g)=(g,0),
\]
since \(\rho(g,0)=\tfrac12\operatorname{Scal}_g\) and \(J(g,0)=0\).

\textbf{Source question.} Which known results concerning nonnegative or positive scalar curvature extend to initial data sets satisfying the strict dominant energy condition? In particular, do the nontrivial homotopy classes detected in the space of positive-scalar-curvature metrics give nontrivial classes in the corresponding space of initial data sets?

A possible precise interpretation of the homotopy question, after fixing the missing hypotheses, is the following: for a specified base metric \(g_0\), degree \(q\geq1\), and specified detected class
\[
0\ne\alpha\in\pi_q(\mathcal R^{>0}(M),g_0),
\]
is
\[
\iota_*(\alpha)\ne0
\quad\text{in}\quad
\pi_q(\mathcal I^{\mathrm{strict}}(M),(g_0,0))?
\]
The source does not identify the classes, detection methods, degrees, or hypotheses under which this implication is intended, nor does it explicitly specify \(\iota\) as the comparison map. It therefore does not justify replacing its question by an unconditional injectivity conjecture. Whether path components are also intended is unspecified. Moreover, a metric with merely nonnegative scalar curvature yields only \(\rho(g,0)\geq|J(g,0)|_g\), so the nonnegative-scalar-curvature portion requires a separate formulation when the target condition is strict.
\subsection{Short English statement}
The question asks whether results about nonnegative or positive scalar curvature, especially detected nontrivial homotopy classes of positive-scalar-curvature metrics, extend to spaces of initial data satisfying the strict dominant energy condition. The manifold hypotheses, topology, comparison map, and classes to be preserved are unspecified, so the displayed homotopy question is a conditional interpretation requiring mathematical review.
\subsection{Sources}
\begin{itemize}
\item Local UnsolvedMath registry, permanent ID \(30004716\), title \emph{Positive Energy Conditions and Spaces of Initial Data}, source problem code \texttt{catalogue entry}; collection: \emph{Oberwolfach Reports Open Problems}. The supplied proposed year is 2021; no individual proposer is supplied.
\item Supplied source citation: \emph{Analysis, Geometry and Topology of Positive Scalar Curvature Metrics} (2022), Oberwolfach Reports, pp.\ 1579--1649, DOI \texttt{10.4171/owr/2021/30}. \url{https://doi.org/10.4171/owr/2021/30}. Also listed among the supplied literature sources as ``Oberwolfach Reports 2022.'' The background describes these reports as documenting mathematical workshops and their open-problem discussions.
\item Supplied literature source: ``Ammann and Gl\"ockle, initial-data index theory.'' \url{https://arxiv.org/abs/2103.11032}. No more specific bibliographic details or theorem locator are supplied.
\item Supplied literature assessment, checked on \texttt{2026-08-22}: the record labels the problem open and states that Dirac--Witten and index-theoretic results treat important spin and strict dominant-energy cases, while no theorem establishes the requested broad generalization for arbitrary initial data sets. This is supplied dated curation derived from the OWR record, not an independently verified literature assessment; its status and cited claims require independent verification.
\item The supplied statement status is \texttt{corrected\_verified}; its accompanying explanation describes a publication-format audit preserving mathematical content while repairing notation or typography. This label does not establish mathematical correctness or independent source verification.
\end{itemize}
\end{document}
